%% pc-onde-periode-longueur — la double périodicité d'une onde progressive.
%% Repère du haut : « photographie » à un instant t fixé -> périodicité spatiale (longueur d'onde).
%% Repère du bas  : enregistrement en un point x fixé   -> périodicité temporelle (période T).
%% Les deux courbes sont identiques : seule la grandeur portée en abscisse change (m / s).
\documentclass[tikz,border=4pt]{standalone}
\usepackage{liaisons}
\begin{document}
\begin{tikzpicture}[x=0.85cm,y=1cm,
  axe/.style={-{Stealth[length=5pt]}, line width=0.6pt},
  courbe/.style={solideA, line width=1.2pt, line join=round, line cap=round},
  cote/.style={{Stealth[length=4pt]}-{Stealth[length=4pt]}, solideD, line width=0.9pt},
  rappel/.style={solideD, line width=0.5pt, dash pattern=on 1.5pt off 1.5pt},
  txt/.style={font=\scriptsize},
  titre/.style={font=\scriptsize\bfseries}]

%% amplitude et sinusoïde commune : 2,5 périodes, période = 3 unités d'abscisse
\def\amp{0.55}
\def\per{3}

%% ================= repère du haut : périodicité spatiale =================
\begin{scope}[shift={(0,4.05)}]
  \node[titre, anchor=south west] at (-0.5,1.28) {Photographie à un instant $t$ fixé};
  \draw[axe] (-0.5,0) -- (8.3,0) node[txt, anchor=north east, inner sep=3pt] {$x$ (m)};
  \draw[axe] (0,-1) -- (0,1.25);
  \node[txt, anchor=south west, inner sep=1pt] at (0.05,1.2) {perturbation};
  \draw[courbe] plot[domain=0:7.5,samples=120,smooth] (\x,{\amp*sin(deg(\x*2*pi/\per))});
  %% cote lambda entre deux crêtes consécutives
  \foreach \c in {0.75,3.75} { \draw[rappel] (\c,\amp) -- (\c,1.02); }
  \draw[cote] (0.75,0.95) -- (3.75,0.95)
      node[midway, fill=white, inner sep=1pt, font=\small, text=solideD] {$\lambda$};
  \node[txt, anchor=north] at (3.75,-0.75) {périodicité \textbf{spatiale}};
\end{scope}

%% ================= repère du bas : périodicité temporelle ================
\begin{scope}[shift={(0,1.2)}]
  \node[titre, anchor=south west] at (-0.5,1.28) {Enregistrement en un point $x$ fixé};
  \draw[axe] (-0.5,0) -- (8.3,0) node[txt, anchor=north east, inner sep=3pt] {$t$ (s)};
  \draw[axe] (0,-1) -- (0,1.25);
  \node[txt, anchor=south west, inner sep=1pt] at (0.05,1.2) {perturbation};
  \draw[courbe] plot[domain=0:7.5,samples=120,smooth] (\x,{\amp*sin(deg(\x*2*pi/\per))});
  %% cote T entre deux maximums consécutifs
  \foreach \c in {0.75,3.75} { \draw[rappel] (\c,\amp) -- (\c,1.02); }
  \draw[cote] (0.75,0.95) -- (3.75,0.95)
      node[midway, fill=white, inner sep=1pt, font=\small, text=solideD] {$T$};
  \node[txt, anchor=north] at (3.75,-0.75) {périodicité \textbf{temporelle}};
\end{scope}

%% ================= relation encadrée =====================================
\node[draw, line width=0.7pt, rounded corners=2pt, inner sep=4pt, align=center,
      font=\small] at (4.2,-0.40)
  {$\lambda = v \times T = \dfrac{v}{f}$\quad
   {\scriptsize (m) \ (m$\cdot$s$^{-1}$) \ (s) \ (Hz)}};
\end{tikzpicture}
\end{document}
